% SOURCE_STATEMENT: "For atoms described by the many-electron Schrödinger equation, determine mathematically meaningful equivalence classes corresponding to chemical groups, determine which atom belongs to which group, and obtain a rigorous electronic-structure description with the simplicity of hydrogen orbital configurations."
% Preparation and final self-audit: CONTRIBUTING.md prompts C and D; source/audit artifacts retained in the local physics_statements_2026_10_11 batch.
% !TeX program = lualatex
\documentclass[11pt]{article}
\usepackage[margin=25mm]{geometry}
\usepackage{fontspec}
\setmainfont{Latin Modern Roman}
\usepackage{amsmath,amssymb,mathtools}
\usepackage{unicode-math}
\setmathfont{Latin Modern Math}
\usepackage{xurl,hyperref}
\hypersetup{colorlinks=true,urlcolor=blue}
\setcounter{secnumdepth}{0}
\setlength{\parindent}{0pt}
\setlength{\parskip}{5pt}
\setlength{\emergencystretch}{3em}
\newcommand{\sref}[2]{\hyperlink{source:#1:#2}{[S#2]}}
\newcommand{\cataloguescope}[1]{\paragraph{Recorded target.}#1}
\newcommand{\R}{\mathbb R}
\newcommand{\C}{\mathbb C}
\newcommand{\N}{\mathbb N}
\newcommand{\Z}{\mathbb Z}
\newcommand{\Q}{\mathbb Q}
\newcommand{\F}{\mathbb F}
\newcommand{\eps}{\varepsilon}
\begin{document}
\section{Mathematical Classification of Atomic Chemical Groups}\label{statement}

\subsection{Definitions and mathematical statement}

\paragraph{Shared notation.}$\mathbb N=\{1,2,\ldots\}$; the standard number fields have their usual meanings. All additional conventions are stated locally.


\paragraph{Definitions and conventions.} For nuclear charge $Z>0$ and $N\ge1$ electrons, the source uses the nonrelativistic atomic Hamiltonian in atomic units,
\[
H_{N,Z}=\sum_{i=1}^N\left(-\tfrac12\Delta_{x_i}-\frac{Z}{|x_i|}\right)+
\sum_{1\le i<j\le N}\frac1{|x_i-x_j|},
\]
on antisymmetric wavefunctions in $L^2((\mathbb R^3\times\{\uparrow,\downarrow\})^N)$. The nucleus is fixed in the Born--Oppenheimer approximation. Total orbital angular momentum, total spin, and parity commute with this Hamiltonian. Their labels on symmetry eigenspaces are $L,S,p$, with eigenvalues $L(L+1)$ and $S(S+1)$ for the squared angular-momentum operators.

\paragraph{Source scope and qualifications.} The source proposes equality of total L,S,p ground-state quantum numbers as a candidate equivalence relation but notes that it need not separate every chemical group. Choosing it is part of the investigation, not a completed derivation of the periodic table. Classification, assigning atoms, and obtaining a simple structural description are all retained.

\paragraph{Catalogue mathematical target.}
\paragraph{Catalogue wording retained for traceability.}
For atoms described by the many-electron Schrödinger equation, determine mathematically meaningful equivalence classes corresponding to chemical groups, determine which atom belongs to which group, and obtain a rigorous electronic-structure description with the simplicity of hydrogen orbital configurations.

\paragraph{Saved literature context.} The supplied catalogue's assessment dated 2026-08-21 records: The passage asks for a rigorous general classification of atomic groups and electronic configurations; large-charge asymptotics answer only restricted ion-sequence cases, not the full question. The assessment is retained as dated provenance; it is not a new comprehensive literature or open-status certification.
\subsection{Short English statement}
Derive atomic groups and a simple, rigorous picture of electronic configurations from the many-electron Schrödinger model.
\subsection{Sources}
\begin{itemize}
\item[\hypertarget{source:statement:1}{S1}] ULAM.AI and the UnsolvedMath Contributors, UnsolvedMath, catalogue entry. Catalogue curation under CC BY 4.0; definitions and English exposition added locally.
\url{https://huggingface.co/datasets/ulamai/UnsolvedMath}.
\item[\hypertarget{source:statement:2}{S2}] Mathematics of Complex Quantum Systems (2009), Oberwolfach Reports, DOI 10.4171/owr/2009/40, pp. 2235--2270; official report source and mathematical context.
\url{https://doi.org/10.4171/owr/2009/40}.
\item[\hypertarget{source:statement:3}{S3}] Official EMS Press PDF of the same Oberwolfach report; PDF pages 16--18, printed pages 2250--2252.
\url{https://ems.press/content/serial-article-files/46241}.
\item[\hypertarget{source:statement:4}{S4}] Friesecke-Goddard large-charge analysis; reference in the saved dated catalogue literature review.
\url{https://doi.org/10.1137/080729050}.
\item[\hypertarget{source:statement:5}{S5}] highly charged ion states; reference in the saved dated catalogue literature review.
\url{https://doi.org/10.1103/PhysRevA.81.032516}.
\end{itemize}
\label{end:statement}
\end{document}
